% Einheit 1 von 5 – Punkte darstellen und Lage lesen (bank/punkte-und-strecken-im-koordinatensystem/e1.jsonl, zusammenbau v0.1)
%% TODO Zeitmarke Einheit 1: Marken-Zeile nicht lesbar
%% TODO Zweigzeile Teil 1: was hier gelernt wird (Satz in Schülersprache aus dem Einheitstitel) – Einheit 1
\einheitenkopf[e1]{Einheit 1 von 5 · Punkte darstellen und Lage lesen}
\zweigzeile{Abitur GK}
\setcounter{aufgabe}{10}

%% TODO Titel: Ich-kann-Satz fehlt in der Bank (Platzhalter: Kettenname) – Nr. 11
%% TODO Anweisung: ein Satz über der ersten Teilaufgabe fehlt in der Bank – Nr. 11
\begin{aufgabe}{Darstellen und Lage lesen}
%% TODO \rechenplatz{n}: Schrittzahl des Grundfalls fehlt in der Bank (nur bei fester Schreibform, 2.3 b)
\begin{teile}
\teil $P(4 | -1 | 0)$ – was sagt die Null? \\ \kreuz{$P$ liegt auf der $x_3$-Achse} \\ \kreuz{$P$ liegt in der $x_1x_2$-Ebene} \\ \kreuz{$P$ liegt in der $x_2x_3$-Ebene}
\teil $A(5 | 3 | 1)$ und $B(2 | 1 | 4)$ eintragen, $C$ ablesen.

\begin{ksys3} \rpunkt{4}{5}{2}{C} \end{ksys3}
C( \leerfeld | \leerfeld | \leerfeld )
\teil Ein Holzwürfel mit Kante $4$ hat eine Ecke im Ursprung und liegt im ersten Oktanten (alle Koordinaten $\ge 0$). Zeichne das Viereck $P(4 | 0 | 1)$, $Q(1 | 4 | 0)$, $R(0 | 4 | 2)$, $S(3 | 0 | 4)$ ein. \hfill (Abitur 2019 GK)

\begin{ksys3} \rquader{0,0,0}{4}{4}{4} \end{ksys3}
\teil Zeichne das Viereck $A(5 | -3 | 0)$, $B(5 | 2 | 0)$, $C(-1 | 4 | 4)$, $D(-1 | -1 | 4)$ in das Koordinatensystem. \hfill (Abitur 2018 LK)

\begin{ksys3} \end{ksys3}
\teil Ein Quader hat eine Ecke im Ursprung $O$ und die Kanten $5$, $6$ und $3$ entlang der $x_1$-, $x_2$- und $x_3$-Achse. $P(5 | 2 | 3)$ liegt auf einer Kante. Zeichne das Dreieck mit den Ecken $O$, $P$ und $R(0 | 6 | 3)$ ein. \hfill (Abitur 2026 GK)

\begin{ksys3} \rquader{0,0,0}{5}{6}{3} \end{ksys3}
\teil Begründe: $P(-3 | 4 | 3)$ und $Q(5 | -1 | 2)$ liegen auf derselben Seite der $x_1x_2$-Ebene. \hfill (Abitur 2025 GK)
\teil Der Grundriss eines Pavillons ist ein Achteck in der $x_1x_2$-Ebene, symmetrisch zur $x_1$- und zur $x_2$-Achse. Gegeben sind $A(6 | 0)$, $B(5 | 3)$, $C(0 | 4)$. Zeichne das vollständige Achteck. \hfill (Abitur 2025 GK)

\begin{ksys}[xmin=-7,xmax=7,ymin=-7,ymax=7,xlabel=$x_1$,ylabel=$x_2$] \punkt{6}{0}{A} \punkt{5}{3}{B} \punkt{0}{4}{C} \end{ksys}
\teil Eine Pyramide hat den quadratischen Boden $ABCD$ mit der Seite $3$ cm; die Spitze $S$ liegt $4$ cm senkrecht über $A$. Im Gitter sind der Boden und die Spitze des Seitendreiecks $ABS$ markiert. Zeichne die drei übrigen Seitendreiecke ein. \hfill (Abitur 2024 GK)

\begin{ksys}[xmin=-5,xmax=9,ymin=-5,ymax=9] \punkt{0}{0}{A} \punkt{3}{0}{B} \punkt{3}{3}{C} \punkt{0}{3}{D} \punkt{0}{-4}{S} \end{ksys}
\teil $B(5 | 1 | 4)$ und $C(1 | 4 | 3)$ sind Ecken des Parallelogramms $ABCD$, $M(2 | 2 | 2)$ ist der Schnittpunkt seiner Diagonalen. Alle Punkte werden parallel zur $x_3$-Achse in die $x_1x_2$-Ebene verschoben. Zeichne $A'B'C'D'$ und $M'$ ein. \hfill (Abitur 2024 LK)

\begin{ksys3} \end{ksys3}
\end{teile}
\end{aufgabe}

%% TODO Titel: Ich-kann-Satz fehlt in der Bank (Platzhalter: Kettenname) – Nr. 12
%% TODO Anweisung: ein Satz über der ersten Teilaufgabe fehlt in der Bank – Nr. 12
\begin{aufgabe}{Darstellen und Lage lesen – Fehler finden, Begründen}
\begin{teile}
\teil Ich finde den Fehler: Lena trägt $P(1 | -4 | 2)$ ein: $1$ in $x_1$-Richtung, $4$ in $x_2$-Richtung, $2$ nach oben.
\teil Begründe: $P(4 | 0 | 3)$ liegt in der $x_1x_3$-Ebene.
\end{teile}
\end{aufgabe}
